% Manuscript skeleton targeting GigaScience / IEEE-ACM TCBB submission format.
% Fill in \author, \affiliation, and body sections as results land.
\documentclass[11pt]{article}

\usepackage[utf8]{inputenc}
\usepackage{amsmath}
\usepackage{graphicx}
\usepackage{booktabs}
\usepackage{hyperref}
\usepackage{algorithm}
\usepackage{algpseudocode}

\title{A Cache-Aligned Sparse Matrix Engine for High-Dimensional \\
       Single-Cell RNA Sequencing Analysis}

\author{%
  % TODO: author list and affiliations
}

\date{}

\begin{document}
\maketitle

\begin{abstract}
Single-cell RNA sequencing (scRNA-seq) datasets routinely exceed one million
cells, and standard sparse-matrix analysis frameworks --- \texttt{scanpy}
(Python) and \texttt{Seurat} (R) --- incur substantial CPU cache thrashing
and out-of-memory failures during $k$-nearest-neighbor (\emph{k}-NN) graph
construction, cell-cell cosine distance computation, and library-size
normalization, owing to the non-contiguous memory access patterns of
Compressed Sparse Row/Column (CSR/CSC) formats. We present
\texttt{Block\_CSR}, a header-only C++20 sparse matrix library built around a
64-byte (cache-line) aligned block-compressed sparse row layout, with
AVX2/AVX-512 SIMD kernels for sparse dot products and cosine similarity, and
zero-copy \texttt{pybind11} interoperability with NumPy/SciPy. % TODO: results summary
\end{abstract}

\section{Introduction}
% TODO: motivate the cache-thrashing / OOM bottleneck (see README.md \S Motivation)
% and position relative to PyNNDescent, Annoy, HNSW.

\section{Related Work}
% TODO: SciPy.sparse.csr_matrix / AnnData memory structures;
% approximate k-NN literature (PyNNDescent, Annoy, HNSW).

\section{Methods}

\subsection{Block\_CSR Data Structure}
% TODO: formal layout spec — alignas(64) value/index arrays, row_ptr semantics,
% relationship to scipy.sparse.csr_matrix. See include/matrix/block_csr.hpp.

\subsection{SIMD Sparse Dot Product and Cosine Similarity}
% TODO: merge-join sparse dot product complexity analysis; AVX2/AVX-512 paths.
% See include/matrix/simd_math.hpp.

\subsection{k-NN Graph Construction}
% TODO: brute-force baseline (include/matrix/knn_graph.hpp) and any
% approximate-index extension.

\subsection{Zero-Copy Python Interoperability}
% TODO: pybind11::array_t buffer-protocol wrapping, capsule-based lifetime
% management (src/python_bindings.cpp).

\section{Benchmarking Protocol}
% TODO: summarize docs/BENCHMARKS.md — datasets, metrics, baselines.

\section{Results}
% TODO: populate once tests/benchmark_1m_cells.py has been run on staged
% 10x Genomics 1.3M Brain Cells data. Target: 20x speedup, 60% memory
% reduction vs. scanpy.pp.neighbors() / scipy.sparse baselines.

\begin{table}[h]
\centering
\caption{Placeholder benchmark results table.}
\begin{tabular}{lccc}
\toprule
Backend & Wall-clock (s) & Peak RSS (GB) & Cache misses \\
\midrule
scanpy.pp.neighbors() & -- & -- & -- \\
scipy.sparse (dense loop) & -- & -- & -- \\
Block\_CSR (this work) & -- & -- & -- \\
\bottomrule
\end{tabular}
\end{table}

\section{Discussion}
% TODO

\section{Reproducibility}
Source code, build instructions, and benchmark scripts are available at the
project repository. % TODO: Compute Capsule / Binder link per \S1.4.

\section*{Data Availability}
% TODO: 10x Genomics dataset citations/DOIs.

\bibliographystyle{plain}
\bibliography{references}

\end{document}
